\documentclass{article}
\usepackage{lmodern}
\usepackage{amsmath}
\usepackage{listings}
\usepackage{fullpage}
\usepackage{parskip}
\usepackage{xcolor}
\lstset{
	basicstyle=\ttfamily,
	columns=fullflexible,
	frame=single,
	breaklines=true,
	postbreak=\mbox{\textcolor{red}{$\hookrightarrow$}\space},
}
\begin{document}
\title{Experiment `p\_6\_2\_b\_shift\_right\_array' Results}
\author{\today}
\date{}
\maketitle
\textbf{Experiment outcome: }SUCCESS
\\\textbf{Bad responses: }0
\\\textbf{Responses containing}~\texttt{assume}~\textbf{: }0
\\\textbf{Responses containing}~\texttt{expect}~\textbf{: }0
\\\textbf{Resolution attempts: }1
\\\textbf{Hard fails (resolution): }0
\\\textbf{Soft fails (resolution): }0
\\\textbf{Verification attempts: }1
\section*{Problem Specification}
\textbf{Problem name: }p\_6\_2\_b\_shift\_right\_array
\\\textbf{Natural language statement: }null
\\\textbf{Method signature: }p\_6\_2\_b\_shift\_right\_array(inputs: array<int>) returns (result: array<int>)
\subsection*{Ensures}
\begin{itemize}
\item \texttt{result.Length == inputs.Length}
\item \texttt{result[0] == inputs[inputs.Length - 1]}
\item \texttt{forall i :: 1 <= i < inputs.Length ==> result[i] == inputs[i - 1]}
\end{itemize}
\subsection*{Requires}
\begin{itemize}
\item \texttt{inputs.Length >= 1}
\end{itemize}
\clearpage
\section*{GenAI interactions}
Below you will find all interactions between the `user' (program) and the `assistant' (OpenAI).
\subsection*{Program $\rightarrow$ GenAI}
\begin{lstlisting}
You are given the following task to perform in Dafny:

Write a method to satisfy the below requirements.

The signature should be:

p_6_2_b_shift_right_array(inputs: array<int>) returns (result: array<int>)

The method should respect the following contract:

requires inputs.Length >= 1, ensures result.Length == inputs.Length, ensures result[0] == inputs[inputs.Length - 1], ensures forall i :: 1 <= i < inputs.Length ==> result[i] == inputs[i - 1]

Produce and show only the Dafny body of this method, including the curly braces that surround it. Do not show the signature nor contract. You must not use 'assume' nor 'expect' anywhere in your code.
\end{lstlisting}
\subsection*{GenAI $\rightarrow$ Program}
\textbf{System fingerprint: }null
\\\textbf{ID: }chatcmpl-DnNikuwRpdYvP0OhCi0eiFi2jBfk1
\\\textbf{Model: }gpt-5.5-2026-04-23
\\\textbf{Created at: }1780661626
\\\textbf{Finish reason: }stop
\\\textbf{Usage: }338 tokens in, and 660 tokens out
\begin{lstlisting}
{
  result := new int[inputs.Length];
  result[0] := inputs[inputs.Length - 1];

  var i := 1;
  while i < inputs.Length
    invariant 1 <= i <= inputs.Length
    invariant result.Length == inputs.Length
    invariant result[0] == inputs[inputs.Length - 1]
    invariant forall j :: 1 <= j < i ==> result[j] == inputs[j - 1]
  {
    result[i] := inputs[i - 1];
    i := i + 1;
  }
}
\end{lstlisting}
\section*{Final Program}
\begin{lstlisting}
method p_6_2_b_shift_right_array(inputs: array<int>) returns (result: array<int>)
	requires inputs.Length >= 1
	ensures result.Length == inputs.Length
	ensures result[0] == inputs[inputs.Length - 1]
	ensures forall i :: 1 <= i < inputs.Length ==> result[i] == inputs[i - 1]
{
  result := new int[inputs.Length];
  result[0] := inputs[inputs.Length - 1];

  var i := 1;
  while i < inputs.Length
    invariant 1 <= i <= inputs.Length
    invariant result.Length == inputs.Length
    invariant result[0] == inputs[inputs.Length - 1]
    invariant forall j :: 1 <= j < i ==> result[j] == inputs[j - 1]
  {
    result[i] := inputs[i - 1];
    i := i + 1;
  }
}
\end{lstlisting}
\section*{Total Token Usage}
\textbf{Input tokens: }338
\\\textbf{Output tokens: }660
\\\textbf{Reasoning tokens: }512
\\\textbf{Sum of `total tokens': }998
\section*{Experiment Timings}
\textbf{Overall Experiment} started at 1780661625727, ended at 1780661673312, lasting 47585ms (47.59 seconds)
\\\textbf{Iteration \#1} started at 1780661625836, ended at 1780661673312, lasting 47476ms (47.48 seconds)
\end{document}
